\exercisesection{Macaulay modules and rings}

\begin{enumerate}
\def\labelenumi{\arabic{enumi})}
\item
  Let A be a Noetherian local ring, and let \(0 \to M' \to M \to M'' \to 0\) be an exact sequence of Macaulay A-modules. Show that we have \(\dim(M') = \dim(M)\) and that \(\dim(M'')\) is equal to \(\dim(M)\) or to \(\dim(M) - 1\) .
\item
  Let A be a Macaulay local ring, and \(\mathfrak{p} \in \operatorname{Spec}(A)\) . Show that we have \(\dim(\widehat{A}/\mathfrak{q}) = \dim(A/\mathfrak{p})\) for every \(\mathfrak{q} \in \operatorname{Ass}(\widehat{A}/\mathfrak{p}\widehat{A})\) . In particular, the ring \(\widehat{\mathrm{A}}/\mathfrak{p}\widehat{A}\) has no embedded prime ideal.
\item
  \begin{enumerate}
  \def\labelenumii{\alph{enumii})}
  \tightlist
  \item
    Let R be a Noetherian, integral, integrally closed Q-algebra, A a ring, and \(\rho: \mathrm{R} \to \mathrm{A}\) an injective homomorphism making A into a finite R-algebra. Prove that the sub-R-module \(\rho(\mathrm{R})\) is a direct summand of A (reduce to cor. 3 of prop. 8).
  \end{enumerate}
\end{enumerate}

\begin{enumerate}
\def\labelenumi{\alph{enumi})}
\setcounter{enumi}{1}
\tightlist
\item
  Let A be a Noetherian local Q-algebra. Prove that A satisfies the following property (the “monomial condition”):
\end{enumerate}

\begin{enumerate}
\def\labelenumi{(\Roman{enumi})}
\setcounter{enumi}{899}
\tightlist
\item
  for every maximal secant sequence \((x_{1},\ldots,x_{d})\) of elements of \(m_{A}\) and every integer p, the element \((x_{1}\ldots x_{d})^{p}\) does not belong to the ideal \((x_{1}^{p+1},\ldots,x_{d}^{p+1})\) . (Reduce to the case where A is complete, hence admits a coefficient field K, and apply a) to a homomorphism \(\rho:K[[X_{1},\ldots,X_{d}]]\to A\) such that \(\rho(X_{i})=x_{i}\) .)
\end{enumerate}

\begin{enumerate}
\def\labelenumi{\arabic{enumi})}
\setcounter{enumi}{3}
\tightlist
\item
  Let A be a local ring, p a prime ideal of A. Let B denote the ring of fractions \(A[Z][S^{-1}]\) , where S is the set of elements \(Z + a\) of A{[}Z{]} with \(a \in \mathfrak{m}_{A} - \mathfrak{p}\) . Let z denote the element Z/1 of B.
\end{enumerate}

\begin{enumerate}
\def\labelenumi{\alph{enumi})}
\item
  Prove that \(\mathrm{B} / (z)\) is isomorphic to \(\mathbf{A}_{\mathfrak{p}}\) and that \(\mathrm{B} / (z - 1)\) is isomorphic to \(\mathbf{A}\) .
\item
  Assume henceforth that the ring A is integral and Noetherian, and that p is generated by a nonzero element p of A. Let C be the ring B{[}T{]}/(z(pT-1)); if A is a Macaulay ring, so is C. Prove that Spec C has two irreducible components \(X_{1}\) and \(X_{2}\) which meet at a closed point x, so that \(\text{codim}(\{x\}, X_{1}) = \text{codim}(\{x\}, X_{2}) = 1\) , \(\text{dim}(X_{1}) = 2\) and \(\text{dim}(X_{2}) \geqslant \text{dim}(A_{\mathfrak{p}})\) . For every integer n, give an example with \(\text{dim}(A_{\mathfrak{p}}) = n\) .
\end{enumerate}

¶ 5) Let A be a Noetherian local ring, M a finitely generated A-module, \(\mathbf{x} = (x_{1},\ldots ,x_{n})\) a sequence of elements of \(\mathfrak{m}_{\mathrm{A}}\) , generating an ideal \(\mathfrak{x}\); assume that the A-module M/xM has finite length (which implies \(n\geqslant \dim (\mathrm{M})\) ). Let \(\mathrm{gr(A)}\) and \(\mathrm{gr(M)}\) denote the associated graded rings of A and M for the x-adic topology, \(\xi_{i}\) ( \(1\leqslant i\leqslant n\) ) the image of \(x_{i}\) in \(\mathfrak{x} / \mathfrak{x}^{2}\) , \(\xi\) the sequence \((\xi_1,\dots ,\xi_n)\) .

\begin{enumerate}
\def\labelenumi{\alph{enumi})}
\item
  For \(p \in \mathbf{Z}\) , we set \(\mathrm{F}^p\mathbf{K}_i(\mathbf{x},\mathrm{M}) = \mathbf{K}_i(\mathbf{x},\mathfrak{x}^{p - i}\mathrm{M})\) (with the convention \(\mathfrak{x}^s = A\) for \(s \leqslant 0\) ). Prove that \(\mathrm{F}^p\mathbf{K}_{\bullet}(\mathbf{x},\mathrm{M})\) is a subcomplex of \(\mathbf{K}_{\bullet}(\mathbf{x},\mathrm{M})\) , and that the complex \(\bigoplus_{n}\mathrm{F}^p\mathbf{K}_{\bullet}(\mathbf{x},\mathrm{M}) / \mathrm{F}^{p + 1}\mathbf{K}_{\bullet}(\mathbf{x},\mathrm{M})\) is identified with \(\mathbf{K}_{\bullet}^{\mathrm{gr}(A)}(\boldsymbol {\xi},\mathrm{gr}(M))\) .
\item
  For every complex of A-modules C such that H(C) is of finite length, we set \(\chi(C) = \sum_{i}(-1)^{i}\lg(\mathrm{H}_{i}(C))\) . Prove that there exists an integer \(p_0\) such that the complex \(\mathrm{F}^{p}\mathbf{K}_{\bullet}(\mathbf{x},\mathrm{M}) / \mathrm{F}^{p + 1}\mathbf{K}_{\bullet}(\mathbf{x},\mathrm{M})\) has zero homology for \(p\geqslant p_0\) ; deduce from this the equalities
\end{enumerate}

\[
\chi \left(\mathbf {K} _ {\bullet} ^ {\operatorname{gr} (\mathrm{A})} (\boldsymbol {\xi}, \operatorname{gr} (\mathrm{M}))\right) = \chi \left(\mathbf {K} _ {\bullet} (\mathrm{x}, \mathrm{M}) / \mathrm{F} ^ {p} \mathbf {K} _ {\bullet} (\mathrm{x}, \mathrm{M})\right) = \sum (- 1) ^ {i} \binom {n} {i} \lg (\mathrm{M} / x ^ {p - i} \mathrm{M})
\]

for \(p \geqslant p_0\) .

\begin{enumerate}
\def\labelenumi{\alph{enumi})}
\setcounter{enumi}{2}
\item
  Prove that the above expression is zero if \(n > \dim(\mathrm{M})\) , and equal to \(e_{\mathfrak{x}}(\mathrm{M})\) if \(n = \dim(\mathrm{M})\) .
\item
  Let \(h \geqslant 0\) , and \(p \geqslant p_{0}\); prove that the canonical injection of \(\mathrm{F}^{p}\mathbf{K}_{\bullet}(\mathbf{x},\mathrm{M})\) into \(\mathrm{F}^{p+h}\mathbf{K}_{\bullet}(\mathbf{x},\mathrm{M})\) is a homology isomorphism. Deduce from this that \(\mathrm{H}_{i}(\mathrm{F}^{p}\mathbf{K}_{\bullet}(\mathbf{x},\mathrm{M}))\) is discrete for the x-adic topology. Prove moreover that the filtration of this module defined by the images of the maps \(\mathrm{H}_{i}(\mathrm{F}^{p+h}\mathbf{K}_{\bullet}(\mathbf{x},\mathrm{M})) \to \mathrm{H}_{i}(\mathrm{F}^{p}\mathbf{K}_{\bullet}(\mathbf{x},\mathrm{M}))\) is x-good, and deduce from this that \(\mathrm{F}^{p}\mathbf{K}_{\bullet}(\mathbf{x},\mathrm{M})\) has zero homology. Conclude that \(\chi(\mathbf{K}_{\bullet}(\mathbf{x},\mathrm{M}))\) is zero if \(n > \dim(\mathrm{M})\) , and equal to \(e_{\mathfrak{x}}(\mathrm{M})\) if \(n = \dim(\mathrm{M})\) .
\end{enumerate}

\begin{enumerate}
\def\labelenumi{\arabic{enumi})}
\setcounter{enumi}{5}
\tightlist
\item
  We keep the hypotheses of the preceding exercise; for every finitely generated A-module N such that N/xN is of finite length, we denote by \(\chi(\mathbf{x},\mathbf{N})\) the integer \(\chi(\mathbf{K}_{\bullet}(\mathbf{x},\mathbf{N}))\) .
\end{enumerate}

\begin{enumerate}
\def\labelenumi{\alph{enumi})}
\item
  For every submodule \(M'\) of \(M\) , one has \(\chi(\mathbf{x}, M) = \chi(\mathbf{x}, M') + \chi(\mathbf{x}, M/M')\) .
\item
  If \(x_{1}\mathrm{M} = 0\), show that \(\chi (\mathbf{x},\mathrm{M})\) is zero.
\item
  Denote by \(\mathbf{x}'\) the sequence \((x_2,\ldots ,x_n)\); prove the equality
\end{enumerate}

\[
\chi (\mathbf {x}, \mathrm{M}) = \chi (\mathbf {x} ^ {\prime}, \mathrm{M} / x _ {1} \mathrm{M}) - \chi (\mathbf {x} ^ {\prime}, \operatorname{Ker} (x _ {1}) _ {\mathrm{M}})
\]

(use exerc. 8 of A, X, p. 207).

\begin{enumerate}
\def\labelenumi{\alph{enumi})}
\setcounter{enumi}{3}
\tightlist
\item
  For \(1 \leqslant i \leqslant n\) , we denote \(\mathrm{K}_i\) the kernel of the homothety with ratio \(x_i\) in \(\mathrm{M}/(x_1\mathrm{M} + \ldots + x_{i-1}\mathrm{M})\) . Prove the equality
\end{enumerate}

\[
\lg (\mathrm{M} / \mathfrak {x} \mathrm{M}) - \chi (\mathbf {x}, \mathrm{M}) = \sum_ {i = 1} ^ {n} \chi ((x _ {i + 1}, \dots , x _ {n}), \mathrm{K} _ {i}).
\]

\begin{enumerate}
\def\labelenumi{\alph{enumi})}
\setcounter{enumi}{4}
\tightlist
\item
  Let \(p_1, \ldots, p_n\) be integers \(> 1\) , \(\mathbf{x}^{\mathbf{p}}\) the sequence \((x_{1}^{p_1}, \ldots, x_n^{p_n})\) . Prove the equality \(\chi(\mathbf{x}^{\mathbf{p}}, \mathrm{M}) = p_1 \ldots p_n \chi(\mathbf{x}, \mathrm{M})\) (first treat the case \(n = 1\) and then reduce to this case using cor. 2 of A, X, p. 157).
\end{enumerate}

¶ 7) Let A be a Noetherian local ring, M a finitely generated A-module. One says that a sequence \((x_1, \ldots, x_n)\) is weakly regular for M if for \(i = 1, \ldots, n\) the kernel of the homothety with ratio \(x_{i+1}\) in \(\mathrm{M}/(x_1\mathrm{M} + \ldots + x_i\mathrm{M})\) is annihilated by \(\mathfrak{m}_A\) .

\begin{enumerate}
\def\labelenumi{\alph{enumi})}
\item
  If \(n \leqslant \dim(M)\), prove that a sequence \((x_{1}, \ldots, x_{n})\) weakly regular for M is secant for M (reason by induction on n). Give an example of a weakly regular sequence that is not secant (one may take \(\mathrm{A} = M = k[X]/(X^{2})\) (where k is a field).
\item
  We say that M is a Buchsbaum module if every secant sequence for M is weakly regular for M. A Macaulay module is a Buchsbaum module; if M is a Buchsbaum module, the \(A_{\mathfrak{p}}\) -module \(M_{\mathfrak{p}}\) is Macaulay for every prime ideal \(p \neq \mathfrak{m}_{\mathrm{A}}\). c) Let M be a Buchsbaum module, \((x_{1}, \ldots, x_{n})\) a secant sequence for M; prove that \(\mathrm{M}/(x_{1}\mathrm{M} + \ldots + x_{n}\mathrm{M})\) is a Buchsbaum module.
\end{enumerate}

\(d\) ) For M to be a Buchsbaum module, it is necessary and sufficient that for every maximal secant sequence \((x_{1},\ldots ,x_{n})\) for M, one has in \(\mathrm{M}_n = \mathrm{M} / (x_1\mathrm{M} + \dots +x_{n - 1}\mathrm{M})\) the equality \(\mathrm{Ker}(x_n)_{\mathrm{M}_n} = \mathrm{Ker}(x_n^2)_{\mathrm{M}_n}\) (first treat the case \(n = 1\) , observing that \(\mathrm{Ker}(x_1)_{\mathrm{M}}\) is then of finite length and reasoning by induction on this length. In the general case, complete every secant sequence \((x_{1},\ldots ,x_{i})\) to a maximal secant sequence \((x_{1},\ldots ,x_{n})\) , and consider the sequences \((x_{1},\ldots ,x_{i - 1},x_{i + 1}^{k},\ldots ,x_{n}^{k},x_{i})\) for \(k\geqslant 1\) ).

¶ 8) Let A be a Noetherian local ring, M a finitely generated A-module. We propose to prove the equivalence of the following two conditions:

\begin{enumerate}
\def\labelenumi{(\roman{enumi})}
\item
  M is a Buchsbaum module;
\item
  there exists a positive integer \(i(M)\) such that one has \(\lg(M/qM)-e_{\mathfrak{q}}(M)=i(M)\) for every ideal \(\mathfrak{q}\subset\mathfrak{m}_{A}\) such that \(M/qM\) is of finite length.
\end{enumerate}

\begin{enumerate}
\def\labelenumi{\alph{enumi})}
\item
  Under hypothesis (i), let x be a secant sequence for M, and let x denote the ideal it generates. With the notations of exerc. 6, d), prove the equality \(\lg(\mathrm{M}/x\mathrm{M}) - e_{\mathfrak{x}}(\mathrm{M}) = \lg(\mathrm{K}_{n})\) . Let \(\mathbf{x}'\) be another secant sequence for M, and let \(\mathrm{K}_{n}^{\prime}\) be the corresponding module; prove the equality \(\lg(\mathrm{K}_{n}) = \lg(\mathrm{K}_{n}^{\prime})\) (reason by induction on \(n\) , considering an element \(t\) of \(\mathfrak{m}_{\mathrm{A}}\) such that the sequences \((x_{1},\ldots,x_{n-1},t)\) and \((x_{1}',\ldots,x_{n-1}',t)\) are secant for M).
\item
  We assume hypothesis (ii) is satisfied. Let \(\mathbf{x} = (x_{1}, \ldots, x_{n})\) be a maximal secant sequence for M, and let \(M_{n} = \mathrm{M}/(x_{1}\mathrm{M} + \ldots + x_{n-1}\mathrm{M})\) ; by applying the formula of exerc. 6, d) to the sequences x and \((x_{1}, \ldots, x_{n-1}, x_{n}^{2})\) (taking account of exerc. 6, e)), prove that one has \(\operatorname{Ker}(x_{n})_{\mathrm{M}_{n}} = \operatorname{Ker}(x_{n}^{2})_{\mathrm{M}_{n}}\) . Conclude with exerc. 7, d)).
\item
  For M to be a Buchsbaum module, it is necessary and sufficient that it be so for \(\widehat{\mathbf{M}}\) ; in this case we have \(i(\widehat{\mathbf{M}}) = i(\mathbf{M})\) .
\item
  Let M be a Buchsbaum module, and \((x_{1},\ldots,x_{p})\) be a sequence of elements of \(\mathfrak{m}_{A}\) secant for M, with \(p<\dim(M)\) ; prove that we have \(i(\mathrm{M}/(x_{1}\mathrm{M}+\ldots+x_{p}\mathrm{M}))=i(\mathrm{M})\) (using the expression of \(i(\mathrm{M})\) given in a)).
\end{enumerate}

\begin{enumerate}
\def\labelenumi{\arabic{enumi})}
\setcounter{enumi}{8}
\tightlist
\item
  Let A be a Noetherian local ring.
\end{enumerate}

\begin{enumerate}
\def\labelenumi{\alph{enumi})}
\item
  Let M be a finitely generated Macaulay A-module, N a submodule of M containing \(\mathfrak{m}_{A}M\). Prove that N is a Buchsbaum module, and that one has \(i(N) = (\dim(M) - 1)[M/N : \kappa_{A}]\) .
\item
  If A is a Macaulay ring, the ideal \(\mathfrak{m}_{A}\) is a Buchsbaum module, which is not Macaulay if \(\dim(A) \geqslant 2\) .
\item
  Let \(k\) be a field; we take for A the subring of \(k[[X,Y]]\) generated by \(X^4, X^3Y, XY^3, Y^4\) . One has \(\dim(A) = 2\) , \(\text{prof}(A) = 1\) , and A is a Buchsbaum ring, with \(i(A) = 1\) (applying \(a\)) by taking for M the integral closure of A).
\end{enumerate}

¶ 10) a) Let \(\rho: A \to B\) be an injective homomorphism of Noetherian rings, such that the sub-A-module \(\mathrm{A}1_{B}\) of B is a direct factor. Prove that if B is a Macaulay ring, then so is A.

(Reduce to the case where A is local and, by reasoning by induction on \(\mathrm{prof}_{\mathrm{A}}(\mathrm{B})\) , where \(\mathrm{Ass}(\mathrm{B})\) contains a maximal ideal. Show, by considering a primary decomposition of 0 in B, that B is then isomorphic to a product, and conclude by reasoning by induction on the number of maximal ideals of B.

\begin{enumerate}
\def\labelenumi{\alph{enumi})}
\setcounter{enumi}{1}
\tightlist
\item
  Let B be a Macaulay ring, G a finite group of automorphisms of B whose order is invertible in B. Prove that the ring of invariants of G is a Macaulay ring.
\end{enumerate}

¶ 11) Let B be a Noetherian integrally closed ring, G a finite group of automorphisms of B, and A the ring of invariants of G.

\begin{enumerate}
\def\labelenumi{\alph{enumi})}
\tightlist
\item
  We assume that for every prime ideal p of \(\mathrm{A}\) of height 1 and every prime ideal q of B lying over p, the valuation \(v_{\mathfrak{q}}\) is unramified with respect to \(v_{\mathfrak{p}}\) (VI, § 8, no. 1). Let \(i: C(A) \to C(B)\) be the canonical homomorphism on divisor class groups (VII, § 1, no. 10); prove that the kernel of \(i\) is isomorphic to the cohomology group \(H^{1}(G,B^{*})\) (A, X, p. 111; let L be the fraction field of B; consider the exact sequence of Z \(^{(G)}\) -modules
\end{enumerate}

\[
0 \rightarrow \mathrm{B} ^ {*} \rightarrow \mathrm{L} ^ {*} \rightarrow \mathrm{D} (\mathrm{B}) \rightarrow \mathrm{C} (\mathrm{B}) \rightarrow 0,
\]

by observing that \(\mathrm{H}^{0}(\mathrm{G},\mathrm{D}(\mathrm{B}))\) is identified with \(\mathrm{D}(\mathrm{A})\) and that \(\mathrm{H}^{1}(\mathrm{G},\mathrm{L}^{*})\) is zero).

\begin{enumerate}
\def\labelenumi{\alph{enumi})}
\setcounter{enumi}{1}
\item
  Let k be a field of characteristic 2; take B to be the ring \(k[X_{1},\ldots,X_{4}]\) , and G to be the group of automorphisms of the k-algebra B generated by the automorphism \(\sigma\) such that \(\sigma(X_{i})=X_{i+1}\) for \(1\leqslant i\leqslant3\) and \(\sigma(X_{4})=X_{1}\) . Deduce from a) that the ring A of invariants of G is factorial.
\item
  Let m be the trace on \(\mathrm{A}\) of the ideal \((\mathrm{X}_1, \ldots, \mathrm{X}_4)\) ; set \(s = \mathrm{X}_1 + \ldots + \mathrm{X}_4\) , \(t = (\mathrm{X}_1 + \mathrm{X}_3)(\mathrm{X}_2 + \mathrm{X}_4)\) , \(u = (\mathrm{X}_1 + \mathrm{X}_2)(\mathrm{X}_3 + \mathrm{X}_4)(\mathrm{X}_1 + \mathrm{X}_4)(\mathrm{X}_2 + \mathrm{X}_3)\) , \(v = \mathrm{X}_1\mathrm{X}_2\mathrm{X}_3\mathrm{X}_4\) . Prove that the sequence \((s, t, u, v)\) is secant but not completely secant for \(\mathrm{A}_{\mathfrak{m}}\) . (Let \(x = \mathrm{X}_1^2(\mathrm{X}_2 + \mathrm{X}_3) + \mathrm{X}_2^2(\mathrm{X}_3 + \mathrm{X}_4) + \mathrm{X}_3^2(\mathrm{X}_4 + \mathrm{X}_1) + \mathrm{X}_4^2(\mathrm{X}_1 + \mathrm{X}_2)\) , \(y = \mathrm{X}_1\mathrm{X}_3 + \mathrm{X}_2\mathrm{X}_4\) , \(w = \mathrm{X}_1\mathrm{X}_2\mathrm{X}_3 + \mathrm{X}_2\mathrm{X}_3\mathrm{X}_4 + \mathrm{X}_3\mathrm{X}_4\mathrm{X}_1 + \mathrm{X}_4\mathrm{X}_1\mathrm{X}_2\) ; prove that we have \(sw + y^2 = u\) and that \(xy + tw\) is divisible by \(s\) , so that \(ux\) belongs to the ideal \((s, t)\) .)
\item
  Show that A is not a Macaulay ring, and that the sub-A-module \(A.1_{B}\) is not a direct summand of B.
\end{enumerate}

